%@ pages 73-86
%@ status translated
% Original section 8, "Localisation de Dwyer-Kan", pp. 73--86.

\clearpage
\section{Dwyer--Kan localisation}
\label{dwyerkanpage}


A good source of examples of simplicial categories, and hence of Segal
$1$-categories, is the theory of simplicial localisations $L(C,W)$
(and their equivalent variants $L^H(C,W)$) developed in the
papers of Dwyer and Kan \cite{DwyerKan1}, \cite{DwyerKan2}, \cite{DwyerKan3}.
If $C$ is a category and $W$ a subcategory one obtains a simplicial
category $L(C,W)$ ``by homotopically inverting the arrows of $W$''
(often the subcategory ``of equivalences'' $W$ is understood, in which
case we shall use the notation $L(C)$, see Convention \ref{convenir}
below). We refer the reader to the works \cite{DwyerKan1},
\cite{DwyerKan2}, \cite{DwyerKan3} for more on the notations and
definitions.


Here the term ``localisation'' does not have exactly
the same meaning as in \S 6 (although there is a relation between the
two notions). For the rest of the paper we shall use the word
``localisation'' to speak of the Dwyer--Kan localisation (or of its
analogue for Segal $n$-categories).

To be historically correct it should be pointed out that Quillen
\cite{Quillen} had already, with
his notion of simplicial closed model category, found
most of the simplicial categories that one constructs with the
Dwyer--Kan localisation; and in particular, in all our examples, the cmc
$M$ that one localises already admits a
simplicial structure, and one could therefore refer to
\cite{Quillen}.
However, the observation of Dwyer--Kan that
the simplicial structure depends,
up to equivalence, only on $M$ and on its subcategory of weak
equivalences is a very important improvement on \cite{Quillen}, which makes it
possible to obtain functorialities easily without having to verify that a
functor preserves the simplicial structure. It is for this reason
that we systematically use
this construction of Dwyer--Kan.


We begin with a more or less trivial remark which, however,
does not appear explicitly in
\cite{DwyerKan1}, \cite{DwyerKan2}, \cite{DwyerKan3}. It is
a variant of Corollary 3.6 of \cite{DwyerKan2} in which one
considers functors that are not necessarily
adjoint and natural transformations of arbitrary direction.
To express this observation, let us agree to call a {\em chain of
natural transformations} between two functors $A$ and $B$ the data
of a sequence $A_0=A, A_1, \dots, A_n=B$ of functors and of a sequence
$u_1, \dots, u_n$ where $u_i$ is a natural transformation
either from $A_{i-1}$ to $A_i$ or from $A_i$ to $A_{i-1}$. Let us also
agree
to say that a natural transformation $u$ between two
functors $F, G: C \rightarrow D$ is with
values
in the subcategory $W$ of $D$ if for every object $X$ of $C$,
$u(X)$ is a morphism of $W$ (likewise for a chain of
natural transformations).

\begin{lemma}
\label{nonadjoint}
Let $C$ and $C'$ be two
categories equipped with subcategories $W$ and $W'$. Let
$$
F: C\rightarrow C' ,\;\;\; G: C' \rightarrow C
$$
be two functors with $F(W)\subset W'$ and $G(W')\subset W$.
Suppose that there exist chains of natural transformations $u:
GF\leftrightarrow
1_{C}$
and $v: FG \leftrightarrow 1_{C'}$
with values in $W$ and $W'$ respectively.
Then $F$ induces an equivalence of simplicial
categories $$
L(C,W)\stackrel{\cong}{\rightarrow}L(C',W').
$$
\end{lemma}
{\em Proof:} The existence of a chain of natural transformations with
values in $W$
implies that the functor induced by
$GF$ on $L(C,W)$ is an equivalence: it then follows from
Propositions 3.3 and
3.5 of \cite{DwyerKan2} that the morphism $L(C,W)(X,Y)\rightarrow
L(C,W)(GFX,GFY)$ is
an equivalence, and it is easy to see that $GF:L(C,W)\rightarrow
L(C,W)$ is
essentially surjective. Likewise the functor induced by $FG$ is an
equivalence. Consequently $F$ and $G$ induce equivalences.
At a stretch, this last step can be seen by using the
closed model category structure for Segal $1$-precategories:
a functor is an equivalence if and only if it induces an
isomorphism in the homotopy category, and, in the latter, one can
use the fact that the existence of right and left inverses implies
invertibility; finally, to conclude, one observes that a functor between
simplicial
categories is an equivalence if and only if the corresponding functor
between Segal $1$-categories is an equivalence.
\eop

The following observation is useful for dealing with subcategories
defined by homotopy-invariant conditions, for
example categories of complexes with cohomology supported in an
interval.

\begin{proposition}
\label{stabilite}
Let $C$ be a category and $W$ a subcategory of $C$ admitting a
homotopy
calculus of fractions
(cf \cite{GabrielZisman}).
Let $B\subset C$ be a full subcategory defined
by a $W$-invariant property, i.e. the inverse image of a subset
$B'$ of the set of equivalence classes of objects of $L(C,W)$.
Then $(B, W\cap B)$ admits a homotopy calculus of fractions and $L(B, W\cap
B)\rightarrow L(C,W)$ is homotopically fully faithful, with as essential
image the full subcategory of $L(C,W)$ that is the inverse image of the
same subset of the set of equivalence classes.
\end{proposition}
This is a direct consequence of the Dwyer--Kan definition
(\cite{DwyerKan2} Definition 6.1).
\eop

The main application of the Dwyer--Kan construction cf \cite{DwyerKan2}
provides, starting from a
closed model category $M$ and from its
subcategory $W$ of weak equivalences, a simplicial
category $L(M,W)$.

\begin{convention}
\label{convenir}
If $M$ is a closed model category
(or a subcategory of fibrant objects, cofibrant objects, etc.) we shall denote
by $L(M)$ the Dwyer--Kan localisation with respect to the subcategory of
weak equivalences. This convention will extend to every category in
which there is a natural notion ``of equivalence''. In case of possible
confusion, e.g. when it is a question of comparing the equivalences for the
coarse topology and the $\Gg$-equivalences, one will put $W$ back into the
notation, for example $L(nSePCh, W^{\rm gro})$ or $L(nSePCh, W^{\Gg})$.
\end{convention}



Let us now point out some methods for computing the ``function complexes'',
i.e. the simplicial sets of morphisms in $L(M)$ where $M$
is a closed model category.
There is of course the direct definition of \cite{DwyerKan1} by means of a
resolution of $M$ by free simplicial categories;
we refer the
reader to \cite{DwyerKan1} for more details.

The second method is that of
``hammocks'', due to Dwyer--Kan \cite{DwyerKan2}. This method is suited to any
localisation but with hammocks of arbitrary length. As is remarked
in \cite{DwyerKan2}, in the case of a closed model category one
can restrict to hammocks of length $3$. For $x,y\in M$ we denote
$$
{\rm ham}^3(M; x,y)
$$
the category of diagrams of the form
$$
x\leftarrow a\rightarrow b \leftarrow y
$$
where the first and the last arrow are weak equivalences
i.e. lie in $W$. The morphisms in ${\rm ham}^3(M; x,y)$
are the diagrams of the form
$$
\begin{array}{ccccccc}
x&\leftarrow &a&\rightarrow &b& \leftarrow &y\\
\downarrow &&\downarrow &&\downarrow && \downarrow \\
x&\leftarrow &a&\rightarrow &b& \leftarrow &y
\end{array}
$$
where the vertical arrows at the two ends are the identities of $x$ and
of $y$.
Dwyer and Kan give in \cite{DwyerKan2} a natural weak equivalence
of simplicial sets between the function complex between $x$ and
$y$, and the nerve of the category of hammocks:
$$
L(M)_{1/}(x,y) \cong \nu  {\rm ham}^3(M; x,y).
$$

The last method is that of ``homotopy function complexes''
(``homotopy function complexes''), due to Reedy \cite{Reedy}, and to Dwyer,
Hirschhorn, Kan \cite{DHK} \cite{Hirschhorn} \cite{DwyerKan3}. This is a
generalisation or weakening of Quillen's method \cite{Quillen}
of considering a simplicial structure on the closed model
category.
In the ``function complexes'' method,
one does not require as much functoriality, which gives
more flexibility (at the price of the fact, as remarked in \cite{DHK}, that the
composition of functions is no longer directly accessible---but this poses
no problem since one already has the category $L(M)$).

Let $x,y\in M$. A {\em fibrant simplicial resolution} of $y$ is a simplicial
object of $M$, ${\bf y}\in M^{\Delta ^o}$, equipped with a morphism
$c^{\ast}y\rightarrow {\bf y}$ where $c^{\ast} y$ is the constant object with
value $y$, such that $y\rightarrow {\bf y}(p)$ is a weak
equivalence for every $p\in \Delta$, and such that ${\bf y}$ is fibrant for the
{\em
Reedy structure} on $M^{\Delta ^o}$ (see \cite{Reedy} \cite{Hirschhorn}
\cite{DHK} \cite{JardineGoerssBook} as well as our discussion further below).
Dually, a {\em cofibrant cosimplicial resolution} of $x$ is a
cosimplicial object of $M$, ${\bf x}\in M^{\Delta}$ equipped with a
morphism ${\bf x}\rightarrow
c^{\ast}x$ such that ${\bf x}(p)\rightarrow x$ is a weak equivalence for
every $p\in \Delta$, and such that ${\bf x}$ is cofibrant for the
Reedy structure on $M^{\Delta}$.

With these resolutions, one obtains (see {\em loc cit.}) a
natural equivalence of simplicial sets
$$
L(M)_{1/}(x,y) \cong d^{\ast} M_{1/}({\bf x}, {\bf y})
$$
where the right-hand term is the diagonal of the bisimplicial set of
morphisms between the two resolutions. If $x$ is cofibrant, one has
$$
L(M)_{1/}(x,y)\cong M_{1/}(x, {\bf y})
$$
while if $y$ is fibrant one has
$$
L(M)_{1/}(x,y)\cong M_{1/}({\bf x}, y).
$$
This is a direct generalisation of the notion of
computing function complexes in homological algebra (by means of
projective and injective resolutions of complexes). This very satisfactory
framework is (to
our knowledge) due to Reedy \cite{Reedy}, and to Dwyer,
Hirschhorn, Kan \cite{DHK} \cite{Hirschhorn} \cite{DwyerKan3}; and it is what
Quillen had in mind when he entitled his book
\cite{Quillen} ``homotopical algebra''. Further on we find the origins in the
theory of
hypercoverings of Verdier \cite{SGA4b}.

If $M$ is a
simplicial closed model category (this is the main case
envisaged by Quillen \cite{Quillen}), one obtains
resolutions from the simplicial structure, hence the localisation
$L(M)$ is
equivalent to
the simplicial category $M^{\rm spl}_{cf}$ of cofibrant
and fibrant objects.

In the introduction of \cite{DwyerKan3}, Dwyer--Kan indicate without
entering into details that $L(M)$ ``captures all the higher homotopy
information implicit in $M$'' that Quillen \cite{Quillen} had
sought, and partially reconstituted in the homotopy category of
$M$---for the non-simplicial case. In fact one can make their assertion more
precise.
Indeed, the following lemma and its corollary allow
the constructions of Gabriel--Zisman that we recalled in \S 7 above
and, in the pointed case, these induce on
$Ho(M)=\tau _{\leq 1}
L(M,W)$
the ``triangulated structure'' defined by Quillen.


\begin{lemma}
\label{prodfibhom}
If $M$ is a closed model category, then $L(M)$ admits
homotopy fibre products and homotopy coproducts in the sense of \S 7
above.
\end{lemma}
{\em Proof:}
We apply the method of resolutions recalled
above. Let
$$
x\rightarrow y\leftarrow z
$$
be a diagram between fibrant objects of $M$, with both arrows fibrant
(every diagram for a fibre product in $L(M)$ is equivalent to a
diagram of this form). For $u\in M$, one chooses a cofibrant
cosimplicial resolution ${\bf u}\rightarrow u$. Then the square
$$
\begin{array}{ccc}
L(M)_{1/}(u, x\times _yz)& \rightarrow & L(M)_{1/}(u, z)\\
\downarrow && \downarrow \\
L(M)_{1/}(u, x)& \rightarrow & L(M)_{1/}(u, y)
\end{array}
$$
is equivalent to the
cartesian square of simplicial sets
$$
\begin{array}{ccc}
M_{1/}({\bf u}, x\times _yz)& \rightarrow & M_{1/}({\bf u}, z)\\
\downarrow && \downarrow \\
M_{1/}({\bf u}, x)& \rightarrow & M_{1/}({\bf u}, y) .
\end{array}
$$
The morphisms at the bottom and on the right are Kan fibrations
(one can verify this by using the fact that ${\bf u}$ is Reedy cofibrant
and that the morphisms $x\rightarrow y$
and $z\rightarrow y$ are fibrant). So this square is also
homotopy cartesian, hence the square of function complexes is
homotopy cartesian. This shows that $x\times_yz$ is a homotopy fibre
product of $x$ and $z$ over $y$.

The proof for homotopy coproducts is dual.
\eop

\begin{corollary}
\label{gzpourM}
If $M$ is a closed model category,
the strict $2$-category $\tau _{\leq 2} L(M)$ satisfies
properties (A), (B), (C) of Gabriel--Zisman and (in the pointed case) their
construction of the operations $\Omega$ and $\Sigma$
applies, and hence gives the ``triangulated structure'' on $Ho\, M=\tau _{\leq
1} L(M)$.
\end{corollary}
The verification is left to the reader.
\eop

For a more general version of this lemma, which allows one to compute the
limits indexed by a small category, see Proposition \ref{calclim}
below (for colimits, see \ref{calccolim}).


We now return to general considerations on the
localisation.
We shall define the localisation for a Segal $n$-category. So let
$A$ be a Segal $n$-precategory and $f\in  A_{1^i}$ an $i$-arrow, $1\leq i\leq
n$. Recall that, with the notations of \cite{limits}, one has
$$
A_{1^i} = Mor _{nSePC}(h(1^i),A)
=Mor _{nSePC}(i\Upsilon (\ast ),A)
$$
where $h(1^i)$ is the Segal $n$-precategory represented by $1^i\in \Theta
^{n+1}$
and where
$i\Upsilon$ denotes the operation $\Upsilon$ of \cite{limits} iterated $i$
times. One has the equality $\Upsilon (\ast )=  I$, where $I$
is the $1$-category having two objects
$0,1$ and one morphism $0\rightarrow 1$.
One may write $i\Upsilon (\ast ) = (i-1)\Upsilon (I)$. Our
$i$-arrow $f$
therefore corresponds to a morphism
$$
f: (i-1)\Upsilon (I)\rightarrow A.
$$
Let $\overline{I}$ be the $1$-category having two
objects $0,1$ and an isomorphism $0\cong 1$. One has $I\subset \overline{I}$,
and hence a cofibration
$$
(i-1)\Upsilon (I)\hookrightarrow
(i-1)\Upsilon (\overline{I})
$$
which one combines with $f$ to define
$$
nSeL(A, f) := A \cup ^{ (i-1)\Upsilon (I)}(i-1)\Upsilon
(\overline{I}).
$$
This is a Segal $n$-precategory (even if $A$ was a Segal $n$-category, it is not
in general a
Segal $n$-category). One must apply the operation $SeCat$ (or a
fibrant replacement) to obtain a Segal $n$-category
$SeCat(nSeL(A, f))$.


If $B$ is a fibrant Segal $n$-category and $u:A\rightarrow B$ a
morphism, then $u$ extends to a morphism
$$
u':  nSeL(A, f) \rightarrow B
$$
if and only if the image $u(f)$ is invertible
up to homotopy in $B$, i.e. if it is invertible as an $i$-arrow
in $\tau
_{\leq i}(B)$.


One can make this universal property a little more precise; but we must
call
on the notations of \S 11 below. The reader is invited
to read \S 11 (see also \cite{nCAT} and
\cite{limits}) before tackling what follows
up to the end of the proof of the next Proposition (marker
$\oslash$).

For $B$ fibrant, the morphism
$$
\underline{Hom}(nSeL(A, f), B)\rightarrow
\underline{Hom}(A, B)
$$
is fully faithful with as essential image the
full sub-Segal-$n$-category of $\underline{Hom}(A, B)$ whose objects
are the $u$ such that $u(f)$ is invertible up to homotopy. This result
(for the proof see the Proposition below) means that
$SeCat(nSeL(A, f))$ possesses the universal property which makes it a
localisation of $A$.

In particular we note that if $f$ was already invertible up to
homotopy in $SeCat(A)$, then the morphism
$A\rightarrow nSeL(A, f)$ is a weak equivalence.

We now extend this construction to the case of
subsets of $i$-arrows. Let $A$ be a Segal
$n$-precategory and let $W=\{ W^i \}$ be a system of subsets of $i$-arrows
$W^i\subset A_{1^i}$, $1\leq i \leq n$. One then defines $nSeL(A,
W)$ as the coproduct of $A$ with a copy of $(i-1)\Upsilon
(\overline{I})$ glued along $(i-1)\Upsilon (I)$, for each $f\in
W^i$. If one wants to have a Segal $n$-category one may then look
at $SeCat(nSeL(A, W))$. One has the same universal property, which we
state in the following form.

\begin{proposition}
\label{propuniSeL}
Let $A$ be a Segal $n$-category and
$W=\{ W^i \}$ a system of subsets of $i$-arrows
$W^i\subset A_{1^i}$. Then $A\rightarrow nSeL(A, W)$ possesses the
following universal
property. If $B$ is a
fibrant Segal $n$-category then
$$
\underline{Hom}(nSeL(A, W), B)\rightarrow \underline{Hom}(A, B)
$$
is fully faithful with as essential image
the full subcategory of the $u: A\rightarrow B$ such that, for every
$f$ in one of the $W^i$,
$u(f)$ is invertible
up to homotopy.
\end{proposition}
{\em Proof:}
In view of the properties of the internal $\underline{Hom}$ cf \ref{interne}
below, it suffices to see this property for the inclusion
$$
(i-1)\Upsilon (I) \rightarrow (i-1)\Upsilon (\overline{I}).
$$
A morphism
$$
(i-1)\Upsilon (I)
$$
is the same thing as a morphism
$$
I\rightarrow B_{1^{i-1}/}.
$$
Since $B_{1^{i-1}/}$ is fibrant, such a morphism extends to $\overline{I}$
if and only if the $i$-arrow in question is invertible up to
equivalence. It remains to prove that the morphism
$$
\underline{Hom}((i-1)\Upsilon (\overline{I}),B)_{1/}(f,g)
\rightarrow
\underline{Hom}((i-1)\Upsilon (I),B)_{1/}(f,g)
$$
is an equivalence; we shall show that it is a trivial fibration by
showing the
lifting property for every trivial cofibration
$E\rightarrow E'$.

For this one must show that the morphism
$$
(i-1)\Upsilon (I) \times \Upsilon (E) \cup ^{
(i-1)\Upsilon (I)\times \{ 0,1\}}
(i-1)\Upsilon (\overline{I})\times \{ 0,1\}
\rightarrow
(i-1)\Upsilon (\overline{I})\times \Upsilon (E)
$$
is a trivial cofibration. The argument is the same as in the
proof of (\cite{limits} Theorem 2.5.1) and we do not reproduce it here.

With this fact, one obtains that for every cofibration $E\rightarrow E'$,
the cofibration
$$
(i-1)\Upsilon (\overline{I})\times \Upsilon (E) \cup ^{
(i-1)\Upsilon (I)\times \Upsilon (E)}
(i-1)\Upsilon (I)\times \Upsilon (E')
$$
$$
\rightarrow (i-1)\Upsilon (\overline{I})\times \Upsilon (E')
$$
% the original prints this formula as two separate displays
is trivial, which gives the desired lifting property
(noting that a morphism
$$
E\rightarrow \underline{Hom}((i-1)\Upsilon (\overline{I}),B)_{1/}(f,g)
$$
identifies with a morphism
$$
(i-1)\Upsilon (\overline{I})\times \Upsilon (E)\rightarrow B
$$
restricting to $f,g$ on $(i-1)\Upsilon (\overline{I})\times \{ 0,1\}$
and likewise for the other morphisms of the square in the lifting
diagram).
\eop


\noindent
$\oslash$

We now return to the case of Segal $1$-categories. We have defined a
localisation $1SeL(A,W)$ by any subset of $1$-arrows $W\subset A_1$.
In this case one does not need the notation $\Upsilon$ since one simply
attaches
copies of $\overline{I}$ along the morphisms $I\rightarrow A$
corresponding to the arrows of $W$. This construction is compatible
with
that of Dwyer--Kan in the following sense.

\begin{proposition}
\label{dklocsegloc}
Let $C$ be a category and $W\subset C$ a subcategory.
Then there is an
equivalence of Segal $1$-categories (compatible with
the identity of
$C$)
$$
L(C,W) \stackrel{\cong}{\rightarrow} 1SeL(C,W)'
$$
where $1SeL(C,W)'$ is a fibrant replacement of
$1SeL(C,W)$.
\end{proposition}
{\em Proof:}
Recall (\cite{DwyerKan1}) that a {\em free category} is a category
freely generated by a set of arrows which we shall call the {\em
generators}. If one fixes a set $O$ of objects, one has the notion of {\em
free product} of a collection of $O$-categories
(i.e. of categories of which $O$ is
the set of objects). If one fixes a set of generators $\{
f_{\alpha}\}$ with $f_{\alpha}$ an arrow from $x_{\alpha}\in O$ to
$y_{\alpha}\in O$, let us denote by $\langle f_{\alpha}\rangle$ the free
$O$-category generated by $f_{\alpha}$. The free $O$-category generated by
$\{ f_{\alpha}\}$ is then the free product of the $O$-categories
$\langle f_{\alpha}\rangle$.

We shall first prove the following result: let $C$ be a free category
(with set $O$ of objects) and
let $x,y\in O$. We shall add an arrow $f$ with source
$x$ and target $y$. Let
$$
C':= C\ast \langle f\rangle
$$
be the free $O$-category generated by the generators of $C$ plus $f$.
On the other hand the arrow $f$ of $C'$ corresponds to a morphism
$I\rightarrow C'$. We shall prove that the morphism of Segal $1$-precategories
$$
C\cup ^{\{ x,y\} }I\rightarrow C'
$$
is a weak equivalence.

To prove this result, we show that $C'$
is obtained from $C\cup ^{\{ x,y\} } I$ by a sequence of explicit
trivial cofibrations. For every $\ell$, let $F_{\ell}C'$ be the simplicial
subset of $C'$ determined by the simplices involving in total, in all
their
principal edges, at most $\ell$ of the generators of $C'$ (i.e. of the
generators of $C$ plus $f$). Here one counts, on one and the same edge (which
corresponds to an arrow of $C'$), the number of generators (with
their multiplicity) that come into play in the arrow in question. Let
$F_{\ell}C$ be the intersection of $F_{\ell}C'$ with $C$, that is to say the
simplicial subset of $C$ of simplices involving at most
$\ell$ generators of $C$ in the above sense. One has the equality
$$
F_1C' = F_1C \cup ^{\{ x,y\} } I.
$$
We shall show that, for $\ell \geq 2$, $F_{\ell}C'$ is obtained
from $F_{\ell -1}C'\cup ^{F_{\ell -1}C}F_{\ell}C$ by coproduct with a
collection of trivial cofibrations. For this, we shall call {\em
elementary $\ell$-simplex} every element $u\in C'_{\ell}$ whose principal
edges are generators of $C'$. Recall also (cf.
\cite{nCAT}) that $h(\ell )$ denotes the simplicial set represented
by $\ell$
(considered as a Segal $1$-precategory constant in the second simplicial
variable).
The morphism
$$
i: h(\ell -1) \cup ^{h(\ell -2)}h(\ell -1)\rightarrow h(\ell )
$$
corresponding to the inclusion of the first and the last face, is a
trivial cofibration of Segal $1$-precategories.
\footnote{
It is easy to see by induction that the morphism
$$
h(m -1) \cup ^{\ast} h(1) \rightarrow h(m )
$$
is a trivial cofibration, and one sees that $i$ is a trivial cofibration
by combining this result for $m=\ell$ and $m=\ell -1$
with the ``three for the price of two'' property.}
If $u$ is an elementary $\ell$-simplex viewed as a
morphism
from $h(\ell )$ to $C'$, one has either $u: h(\ell )\rightarrow C$, or
$$
u^{-1} (F_{\ell -1}C' \cup ^{F_{\ell -1}C} C) =
i\left( h(\ell -1) \cup ^{h(\ell -2)}h(\ell -1)\right) .
$$
In the first case one does not need to add $u$; in the second case one
can add to $F_{\ell -1} C'\cup ^{F_{\ell -1}C} F_{\ell}C$ a copy of
$h(\ell )$ glued along $ h(\ell -1) \cup ^{h(\ell -2)}h(\ell -1)$,
corresponding to $u$. Let us provisionally call $B$ the result of all these
coproducts, one for each elementary $\ell$-simplex not contained
in $C$. There is an obvious morphism $B\rightarrow C'$ and it is easy
to see
that this morphism is injective, because an element (i.e. $k$-simplex)
introduced
by the cofibration associated with $u$ knows all the generators
appearing
in $u$, in particular it cannot come from a
cofibration associated with a $u'\neq u$. On the other hand the image of this
morphism is
visibly $F_{\ell} C'$. So $B=F_{\ell} C'$ and we have proved that
$$
F_{\ell -1} C'\cup ^{F_{\ell -1}C} F_{\ell}C\rightarrow
F_{\ell} C'
$$
is a trivial cofibration. The argument also shows that
$$
F_{\ell -1} C'\cup ^{F_{\ell -1}C} C\rightarrow
F_{\ell } C'\cup ^{F_{\ell }C} C
$$
is a trivial cofibration. The composition of these
trivial cofibrations gives the trivial cofibration
$$
F_1C'\cup ^{F_1C} C\rightarrow C',
$$
that is to say the result sought.


Now we continue the proof of the Proposition. Let us denote by
$A\ast B$ the free product of two $O$-categories, and by $A\cup ^OB$ the
coproduct as $O$-precategories of Segal (i.e. Segal $1$-precategories having $O$
as set of objects). The result above implies, by
induction (plus
possibly passage to a filtered colimit, which poses no problem), that if $C$
is the free category generated by $\{
f_{\alpha}\}$
then
$$
\coprod _{\alpha}^{O} \langle f_{\alpha} \rangle \rightarrow C
$$
is a weak equivalence of Segal $1$-precategories. It follows that if
$A$ and
$B$ are free $O$-categories, then the morphism
$$
A\cup ^OB\rightarrow A\ast B
$$
is a weak equivalence of Segal $1$-precategories.

If $A$ is a simplicial $O$-category one may consider it as a
Segal $1$-category (using the simplicial variable as second
simplicial variable for the Segal category). The result above
implies directly that if $A$ and $B$ are simplicial $O$-categories,
free at each stage, then the morphism
$$
A\cup ^OB\rightarrow A\ast B
$$
is a weak equivalence. Now, the coproduct (in a cmc) of
cofibrant objects (and for \ref{SeCmf} every object is cofibrant)
preserves the equivalence type. Dwyer--Kan show that the free product of
arbitrary simplicial
$O$-categories preserves equivalence
(\cite{DwyerKan1} Proposition 2.7). Using free simplicial
resolutions (\cite{DwyerKan1}) one obtains that for every pair of simplicial
categories $A,B$ having the set $O$ as objects,
the morphism
$$
A\cup ^OB\rightarrow A\ast B
$$
is a weak equivalence of Segal $1$-precategories.

We now return to the definition of the Dwyer--Kan localisation
\cite{DwyerKan1}. If $C$ is a category and $W$ a subcategory,
Dwyer--Kan choose a free resolution of $C$ of the form
$A\ast B$ with $A$ and $B$ simplicial categories free at
each level
(and sharing the same set of objects $O$),
and such that $B$ is a resolution of $W$. Let us denote by $\overline{B}$
the simplicial category obtained by inverting the arrows of each level
of $B$. By definition one has
$$
L(C,W):= A\ast \overline{B}
$$
while
$$
1SeL(A\ast B, B)= (A\ast B)\cup ^{B}1SeL(B,B).
$$
Observe moreover that $1SeL(B,B)\rightarrow \overline{B}$ is a weak
equivalence (for this one may suppose that $B$ is a free category, in which
case
$\overline{B}$ is its $0$-groupic completion and $1SeL(B,B)$ as well).

One
has a sequence of arrows $$
A\cup ^O \overline{B} = (A\cup ^OB)\cup ^B\overline{B}\rightarrow
(A\ast B)\cup ^{B}\overline{B}\rightarrow A\ast \overline{B},
$$
and by the remarks above,
the first arrow and the composite are
weak equivalences of Segal $1$-precategories. By
``three for the price of
two'' one obtains that the arrow
$$
1SeL(A\ast B, B)= (A\ast B)\cup ^{B}\overline{B}\rightarrow A\ast
\overline{B}=L(C,W)
$$
is an equivalence. On the other hand the localisation $1SeL$
preserves equivalences, so the arrow
$$
1SeL(A\ast B, B)\rightarrow 1SeL(C,W)
$$
is a weak equivalence; hence, choosing a fibrant replacement
$1SeL(L,W)'$ % the French prints $1SeL(L,W)'$ here
one obtains an equivalence (essentially well defined)
of Segal $1$-categories
$$
L(C,W)\stackrel{\cong}{\rightarrow} 1SeL(C,W)'.
$$
This completes the proof.
\eop


\begin{corollary}
\label{dklocseglocspl}
Let $C$ be a simplicial category and $W\subset C$ a simplicial subcategory.
One may suppose that $W$ is saturated in the sense that $W(x,y)$ is
a union of connected components of $C(x,y)$ (cf \cite{DwyerKan1}).
The set $W_{1,0}$ of vertices
of the simplicial set of arrows of $W$ is a subset of
$1$-arrows of the Segal $1$-category corresponding to $C$. One then has
an
equivalence of Segal $1$-categories (compatible with the injections of
$C$)
$$
L(C,W) \stackrel{\cong}{\rightarrow} 1SeL(C,W_{1,0})'.
$$
\end{corollary}
{\em Proof:}
One supposes $C$ free as in \cite{DwyerKan1}, and one applies the
Proposition
at each level. One may note that by \cite{DwyerKan1}, $L(C,W)$ is
equivalent to $L(C, W_{1,0})$.
\eop

\begin{corollary}
\label{propunivM}
Let $C$ be a simplicial category and $W\subset C$ a simplicial subcategory.
If $B$ is a fibrant Segal $1$-category, then
the morphism of Segal $1$-categories
$$
\underline{Hom}(L(C,W), B)\rightarrow \underline{Hom}(C,B)
$$
is fully faithful with as essential image the
subcategory of the morphisms $C\rightarrow B$ which send
the elements of $W$ to
invertible morphisms (up to equivalence) in $B$.
\end{corollary}
{\em Proof:}
It suffices
to apply Proposition \ref{propuniSeL} and \ref{dklocseglocspl} (or just
\ref{dklocsegloc} if $C$ is a non-simplicial $1$-category).
\eop

This corollary should have a variant formulated in terms of
simplicial categories with the simplicial categories
$Coh(\cdot , \cdot )$ of morphisms of Cordier and Porter, see the problems
\ref{probleme1}--\ref{probleme2} above.

{\em Remark:} Strictly speaking, there does not necessarily exist a morphism
$C\rightarrow L(C,W)$. On the other hand, if $\widehat{C}\rightarrow C$
is the standard
simplicial resolution of $C$ (which is an equivalence of simplicial
categories) then by construction \cite{DwyerKan1} one has a morphism
$\widehat{C}\rightarrow L(C,W)$. In the rest of the paper we shall
neglect
this technical point (we have already ignored it in the statement of the
preceding corollary) and {\em pretend} that one has a morphism
$C\rightarrow L(C,W)$.
This difficulty does not concern $1SeL(C,W)$,
which by construction receives
a morphism with source $C$.

\bigskip



The following corollary is the analogue, for Segal categories with
several objects, of a fact well known for topological monoids: every weak
monoid (i.e. object of the category underlying
a delooping machine; for example a Segal category with a single
object, for the Segal machine) is equivalent to a strict topological
monoid. See the work of
Fiedorowicz \cite{Fiedorowicz}, or Proposition 1.12 of Dunn
\cite{Dunn}.

%%% added in v2:
In fact this corollary is to be found
in \cite{DKS}, which contains for the first time
the notion of Segal $1$-category (the present remark
was added in version 2).


\begin{corollary}
\label{catsimpl}
If $A$ is a fibrant Segal $1$-category then there exists a simplicial
category $C$ and an equivalence $C\rightarrow A$.
\end{corollary}
{\em Sketch of proof:}
By considering $A$ as a simplicial presheaf on the category
$\Delta$, one may apply the ``barycentric subdivision''
above each object of $\Delta$ (see Lemma \ref{bd} below) so as
to obtain a presheaf of categories on $\Delta$. Then one
applies the Grothendieck integration operation (cf \S 16 below), by
using an analogue of Lemma \ref{bd} for presheaves of categories
above $\Delta$ instead of simplicial sets. This is the analogue
for Segal categories of
the platification (``flattening'') used in
Dwyer--Kan (\cite{DwyerKanDiags} ``A delocalization theorem'' 2.5). One thus
obtains a $1$-category $\beta (A)$ equipped with a subset $\delta (A)$ of
$1$-arrows: one includes in it all the ``vertical'' arrows, i.e. those
of the fibre categories, as well as the ``horizontal'' arrows which
enter into the $\delta (A)$ of \ref{bd}. As in \ref{bd},
one has
$$
A\cong 1SeL(\beta (A), \delta (A))'
$$
(equivalence in the homotopy category of $1SeCat$).
Then, by the preceding Proposition, one also has
$A \cong L(\beta (A), \delta (A))$. And since $A$ is fibrant,
this equivalence can be realised as an equivalence
$L(\beta (A), \delta (A))\rightarrow A$.
\eop



This proof also shows that there exists a simplicial set $X$
satisfying
$$
SeCat(X) \cong A.
$$
Indeed, $\beta (A)$ is a $1$-category which
corresponds to a simplicial set (its nerve); and the localisation
$1SeL(\beta (A), W)$ is obtained by gluing copies of
$\overline{I}$ along the morphisms $I\rightarrow \beta (A)$ of $W$.
The coproduct stays within simplicial sets (indeed, the
result stays
constant in the second simplicial variable of the Segal $1$-precategories).
One therefore indeed has a simplicial set
$$
X:= 1SeL(\beta (A), W)
$$
satisfying $A\cong SeCat(X)$.


\subnumero{Homotopy adjoints}

The property of adjunction is transmitted to Dwyer--Kan localisations
and becomes a property of homotopy adjunction. Let $A$ and $B$ be two
Segal $n$-categories (which one may suppose fibrant), and
$$
F: A\rightarrow B,\;\;\;\; G: B\rightarrow A
$$
two morphisms. Let also
$$
u\in \underline{Hom}(A,A)_{1/}(1_A, GF)
$$
be a natural transformation. Then
for every $x\in A_0$ and every $y\in B_0$, $u$ induces a morphism
(essentially
well defined)
$$
B_{1/}(Fx,y)\rightarrow A_{1/}(x,Gy).
$$
We shall say that $F$ and $G$ are {\em homotopically adjoint via $u$} if
this morphism is an equivalence of Segal $n-1$-categories for every
$x$ and every $y$. We shall say that $F$ is {\em the left adjoint} and $G$ {\em
the right
adjoint}. One could also construct the definition starting from
a natural
transformation
$$
v\in \underline{Hom}(B,B)_{1/}(FG,1_B).
$$
And to do things properly, one should verify
that the two definitions coincide, that the adjoint
$(G,u)$ of a morphism $F$ is essentially unique if it exists, etc.

%%% sentence added in v2 and v3:
See Cordier--Porter \cite{CordierPorter} for the case $n=1$.

Likewise one could envisage defining the notion of an adjoint pair of
$1$-arrows in an arbitrary Segal $n$-category (and recover the
notion above with the $n+1$-category $nSeCAT$). This could correspond
to what
Baez--Dolan call a ``dual'' in an $n$-category
\cite{BaezDolan}.

\begin{lemma}
\label{adjointh}
Let $F:M\rightarrow N$ be a left Quillen functor between closed model
categories, and let $G$ be its right adjoint.
Then
$$
L(F): L(M)\leftrightarrow L(N): L(G)
$$
are homotopically adjoint.
\end{lemma}
{\em Proof:}
First one must note that the computation of the $L(M)_{1/}(x,y)$ by
resolutions is
compatible with composition as far as possible: if for example
$y\rightarrow {\bf y}$ is a fibrant simplicial resolution, and if
$x\rightarrow x'$ is a morphism between cofibrant objects, then
the induced morphism
$$
M_{1/}(x',{\bf y})\rightarrow M_{1/}(x,{\bf y})
$$
coincides with the composition morphism
$$
L(M)_{1/}(x',y)\rightarrow L(M)_{1/}(x,y)
$$
via the equivalence of \cite{DwyerKan2} which was recalled above.

We apply this to the adjunction morphism $x\rightarrow GFx$.
One obtains that if $x$ is cofibrant and if $y\rightarrow {\bf y}$ is a
fibrant simplicial resolution, then the morphism
$$
L(N)_{1/}(Fx,y)\rightarrow L(M)_{1/}(x,Gy)
$$
induced by the adjunction morphism coincides (up to equivalence)
with the morphism
$$
N_{1/}(Fx, {\bf y})\rightarrow M_{1/}(x,G{\bf y})
$$
(here the fact that $F$ and $G$ are Quillen functors (left and
right) implies that $Fx$ is cofibrant, and that $Gy\rightarrow G{\bf y}$ is a
fibrant simplicial resolution). Now this last morphism is an isomorphism,
so the adjunction morphism on the localisations is an equivalence.
\eop



For another property of a closed model category
$M$ which can
be ``amplified''
into a property of $L(M)$ see Proposition
\ref{hominterne} on homotopy internal $Hom$'s.
We conclude by posing the analogous problem for ``derived
functors''.

\begin{problem}
\label{derive}
Define the notion of {\em homotopy derived functor} in the framework of
simplicial categories or Segal $1$-categories (or possibly Segal
$n$-categories). Prove that if $M$ and $N$ are closed model
categories and if $F: M\rightarrow N$ is a Quillen functor (left,
say), then the composite functor
$$
L(M)\cong L(M_c)\rightarrow L(N_c)\cong L(N)
$$
is the homotopy derived functor of $M\rightarrow L(N)$ along
$M\rightarrow L(M)$. Give a statement of essential uniqueness for the
homotopy derived functor, if it exists.
\end{problem}

The solution of
this problem would make it possible to see that the Segal $1$-prestack
$L(\mm )$ associated with a left Quillen presheaf $\mm$ (see \S 17)
is essentially well defined independently of the choice of the
subcategories
of cofibrations in the $\mm (y)$.
